\chapter{Conclusões e Perspectivas}

Esta dissertação teve como objetivo estabelecer a ligação entre a formulação lagrangiana e as simetrias de uma teoria de campos, a construção do tensor energia--momento e da energia do vácuo, e a emergência do efeito Casimir quando o espectro de modos do campo é modificado pela presença de fronteiras. Ao longo do trabalho, esse programa foi desenvolvido em dois níveis complementares: primeiro, no contexto fundamental de um campo escalar submetido a condições de contorno em placas paralelas; em seguida, em uma extensão do problema para um cenário com dimensão extra compactificada e acoplamento éter--escalar.

O ponto de partida foi a dinâmica clássica de um campo escalar real descrita pela equação de Klein--Gordon \eqref{eq:KLeinGordonEq} e sua generalização na presença de uma fonte externa \eqref{eq:GenKleinGordonEq}. A partir desse formalismo, evidenciou-se que a informação essencial para a discussão do vácuo e de seus observáveis encontra-se na definição da ação e da densidade lagrangiana, Eqs.~\eqref{eq:acao} e \eqref{eq:lagphi}, e nas consequências das simetrias contínuas dessa ação.

Em particular, explorou-se a condição de invariância da ação sob transformações gerais \eqref{eq:tranformacoes}, com ênfase no caso em que a variação da lagrangiana se reduz a uma divergência total \eqref{eq:deltaL1}. A equivalência entre essa condição e a existência de uma corrente conservada foi explicitada ao comparar \eqref{eq:deltaL1} com a identidade obtida via Euler--Lagrange \eqref{eq:deltaL2}.

Como exemplo central, analisou-se a simetria de translação no espaço-tempo \eqref{eq:spacetimetranslation}, que induz a transformação infinitesimal do campo \eqref{eq:deltaPhi=epsilonPartialPhi} e implica a forma de divergência total \eqref{eq:partialLPartialK}. Nesse contexto, identificou-se o tensor energia--momento como a corrente de Nöther associada à translação, Eq.~\eqref{eq:stressTensor}, e sua conexão com o operador energia--momento é obtida pela integração espacial (ver, por exemplo, a relação final em \eqref{eq:HeP}). Essa etapa é conceitualmente decisiva, pois o efeito Casimir é extraído de valores esperados do tensor energia--momento e de sua dependência geométrica.

A transição para a teoria quântica foi realizada via quantização canônica e expansão em modos. Para contornos estáticos, introduziram-se soluções de frequência positiva e negativa \eqref{eq:campos} e o correspondente problema espectral espacial \eqref{eq:spatialEq}, com produto escalar \eqref{eq:scalarProduct} e condição de normalização \eqref{eq:normalization}. O operador de campo \eqref{eq:campogeral} foi construído em termos de operadores de criação e aniquilação que satisfazem as relações de comutação \eqref{eq:commutrel}, com definição do vácuo \eqref{eq:aKet0=0}.

Nesse formalismo, a energia de vácuo surge como valor esperado da componente $00$ do tensor energia--momento \eqref{eq:tensorMeanvalue}. Substituindo-se a solução geral \eqref{eq:campogeral} em \eqref{eq:tensorMeanvalue}, obtém-se a expressão local \eqref{eq:<T>}, cuja integração no volume leva à soma sobre modos para a energia de ponto zero \eqref{eq:modeSum}. Em contrapartida, na ausência de fronteiras, a solução plana \eqref{eq:phiLivre} conduz à densidade de energia de vácuo de Minkowski \eqref{eq:freeSol}. Embora cada energia de vácuo seja divergente, as diferenças entre configurações com e sem restrições podem ser finitas e mensuráveis.

No Capítulo~2, o formalismo acima foi aplicado à geometria de duas placas paralelas com condições de contorno de Dirichlet. O espectro do campo é explicitamente modificado pela discretização do número de onda na direção normal às placas, o que se reflete na densidade de energia de vácuo entre as placas, Eq.~\eqref{eq:comPlacas}. A densidade correspondente do vácuo de Minkowski, escrita como integral contínua, é dada por \eqref{eq:semPlacas}. A densidade renormalizada \eqref{eq:denEnPlacas} foi obtida a partir da diferença de energia nas configurações com e sem placas usando Abel--Plana \eqref{eq:abelPlana}.

Ao descartar o termo oscilatório de \eqref{eq:denEnPlacas}, que não contribui para a energia total, obteve-se a energia de Casimir por unidade de área no limite não-massivo, Eq.~\eqref{eq:EnergiaEntrePlacas}. Finalmente, a pressão de Casimir foi extraída como derivada da energia em relação à separação entre as placas, culminando no resultado padrão \eqref{eq:pressaoEntreAsPlacas}. Esses resultados explicitam que o efeito Casimir é, essencialmente, um efeito espectral: as fronteiras alteram o conjunto de frequências \eqref{eq:campos} que contribuem na soma \eqref{eq:modeSum}, e a parte finita dessa reorganização espectral se manifesta como uma força (ou pressão) dependente de $a$.

No Capítulo~3, essa estrutura foi generalizada para o caso de um campo escalar acoplado a um campo éter em um espaço-tempo com topologia $R^4 \times S^1$. A compactificação da quinta dimensão e a imposição de condições de quase-periodicidade introduziram uma nova escala no problema e modificaram a relação de dispersão dos modos do campo. Como consequência, a energia de Casimir e a pressão entre as placas passam a depender não apenas da separação $a$, mas também do raio de compactificação, do parâmetro de quase-periodicidade e da razão associada ao acoplamento éter--escalar.

A análise desenvolvida nesse capítulo mostrou que o campo éter atua efetivamente como um fator de redução de comprimento, alterando a escala associada à dimensão extra e influenciando o comportamento da energia e da força de Casimir. Além disso, a obtenção da forma regularizada da energia permitiu discutir regimes específicos dos parâmetros do modelo e identificar como a presença simultânea da dimensão extra compactificada e do acoplamento LIV modifica o espectro de flutuações do vácuo.

\section{Perspectivas e trabalhos futuros}
\label{sec:perspectivas}

A principal continuação natural deste trabalho consiste em refinar o cenário fenomenológico estudado no Capítulo~3. Em particular, a comparação com os dados experimentais pode ser aprimorada pela incorporação de efeitos que não foram incluídos explicitamente no modelo, como correções de temperatura finita, condutividade finita dos materiais e outras contribuições associadas ao aparato experimental. Esse refinamento é importante para avaliar com maior precisão o alcance das restrições obtidas para os parâmetros do modelo e para o tamanho da dimensão extra.

Outra direção promissora é a inclusão de campos externos. Em especial, a introdução de um campo magnético externo exige a substituição da derivada ordinária por uma derivada covariante, modificando a dinâmica do campo escalar, o problema espectral associado e, consequentemente, a estrutura da energia de vácuo e da força de Casimir. Essa generalização permite investigar a competição entre efeitos geométricos, topológicos e de acoplamento externo na organização do espectro de modos.

Também permanecem em aberto extensões para geometrias mais gerais e para outras classes de condições de contorno, bem como uma análise mais detalhada de limites particulares do modelo, como o regime em que o parâmetro de quase-periodicidade tende a zero. Em todos esses casos, o núcleo metodológico permanece o mesmo: determinar o espectro permitido, construir a energia de vácuo renormalizada e extrair as grandezas observáveis associadas.

Em síntese, os resultados desta dissertação mostram que o efeito Casimir fornece um teste conceitual e técnico privilegiado para a teoria quântica de campos. A partir do formalismo construído nos capítulos iniciais e de sua aplicação a um cenário com dimensão extra compactificada e acoplamento éter--escalar, tornou-se possível compreender como informações geométricas, topológicas e dinâmicas entram no espectro do campo e se traduzem em grandezas observáveis, como energia, pressão e deslocamentos de frequência. Assim, o trabalho estabelece uma base sólida tanto para o estudo do efeito Casimir em geometrias confinadas quanto para investigações futuras envolvendo extensões da estrutura do espaço-tempo e quebras efetivas de simetria de Lorentz.
